\section{Evaluation}
\subsection{Simulation Testing}
The best way to test the accuracy of the finished product is in the desired market location (in a vehicle). To collect the data, I
established a simple program that will capture 5-minute videos of the course of an hour while you drive. Additionally, other datasets,
such as the SUST \cite{SUST} and BINA \cite{BINA}will be used to help increase the variety and test range for the model.

The duration of the videos is not importaint, but I tough it was a good idea to break an extended recording session down into sub-sessions,
incase removing power mid recording resulted in its corruption. I was hoping this would not be an issue (which it wasn't), but based on my
past experiences with varying dashcameras, it was a risk I was not willing to take. From here, frames were extracted at regular intervals.
These were then manually classified based on my own opinion of tiredness (the same opinion that the model was trained on). Then these
frames were passed through the model. All of this data was then combined into a single datasets with the following attributes:
\begin{itemize}
    \item \textbf{Actual classification} This is the classification I assigned to the image based on my own opinion of tiredness.
    \item \textbf{Predicted classification} This is the classification that the model assigned to the image.
    \item \textbf{Darkness} This is a value between 0 and 1 that represents how dark the image is. This was calculated by taking the average
    pixel value across the entire image and dividing it by 255 (the maximum pixel value).
    \item \textbf{Gender} This is the gender of the driver in the image.
    \item \textbf{Camera position} The position of the camera relative to the driving position.
    \item \textbf{Accessories} This is a binary value that represents whether the driver is wearing any accessories (e.g. glasses, hats,
    masks).
    \item \textbf{Hair in face} A binary value that represents whether the driver has any hair in their face (e.g. fringes, beards).
    \item \textbf{Alert\_Confidence} This is a direct output from the model and represents how sure the model is that they are
    \textbf{alert} (not tired)
    \item \textbf{Tired\_Confidence} This is a direct output from the model and represents how sure the model is that they are
    \textbf{tired}.
\end{itemize}
This will then be used to form a basis for my \textbf{Stratified Stress Testing}.
Overall, I would like to see how each factor effects the model as well as how well the model general performs in everyday simulations.


\subsection{Data Extraction:}
\paragraph{How many of each classification?}
To determine if there is a trend between the predicted result and that attribute(s) it should fulfill the following:
\begin{equation}
    E = z \cdot \sqrt{\frac{p(1 - p)}{n}}
\end{equation}
To achieve the ‘scientific gold standard’ of a 95\% confidence that a pattern exists within the data, I have taken p=0.5
(the worst-case probability, equal to 50/50 a guess), n=number of images and E being the margin or error. Given this you need 384
images to provide a confidence of 95\% within the data, and so, to cover all bases, I chose to aim for at least 400 images of each
attribute (at a minimum). This can help reduce Type II errors by ensuring that there is sufficient evidence of this hypothesis.

\paragraph{How data was extracted?}
While simple and not necessarily relevant for the overall project, this was formed using the ****** library. This searched through a
selected folder and displayed every image within it; any videos were dissected into frames (captured 5 seconds apart). From the videos,
the frames were stored in temporary files to allow me to navigate through while removing the delay of having to search through a video
for an image every time I wanted to move on (or to previous frames in the event of misclassification). Once this has been done, I manually
went through every image and applied labels to each feature. While certain aspects such as darkness would be done automatically,
determining the camera position relative to the driver would have been more challenging (simply taking camera position to relative head
direction, would not be as informative as position as for an articular user, this is fixed following setup).
\begin{figure}[H]
    \centering
    \includegraphics[width=0.8\textwidth]{images/data_extraction.png}
    \caption{The interface for data extraction. Features an image from the SUST dataset \cite{SUST}}
\end{figure}
All this combined into a 5676-line CSV file which would hopefully provide an insight into positive or negative trends within the model.

\subsection{Preliminary Observations}
While assigning variables, I noticed several trends in the data that I would like to point out prior to the quantitative data.

With the camera position being in a high-up position, the data tended to miss classify more than in other positions. I believe this to be
because of how your eyes are orientated. With the camera being in a higher position, the eyelids are perceived as being more closed.
This bias could potentially be represented in the data.

Eye shape. Some of the data contained people whose eyes are considered more of a ‘hunter’ eye. For these people, it seemed that the miss
classification of tiredness was more prevalent that other people. This and other factors also appear to lead to one person having a
higher tendency to one classification over another.

Also, for those who would sit more upright (or laid back) in their seats also seemed to be mislabeled more. This also seemed true for
incredibly sunny environments for the same reason; they wouldn’t need their eyes as open to be able to see and so were considered tired
incorrectly.

Another strange observation was that when someone was wearing a mask (even just under the chin) the person was often considered to be
tired incorrectly, compared to other accessory features.

\subsection{Quantitative Analysis}
For this, I am using Orange \cite{orange}. This is a datamining software that helps identify features in both a
visual and numerical way. This is used in industry for rapid prototyping, data visualization and exploratory data analysis.

\subsubsection{General Data Spread:}
For this I used a ‘column statistics’ node to express how each attribute is distributed over the entire dataset.
\begin{figure}[H]
    \centering
    \includegraphics[width=0.6\textwidth]{images/data_spread_nodes.png}
    \caption{Node structure for general data spread}
\end{figure}
\begin{figure}[H]
    \centering
    \includegraphics[width=0.9\textwidth]{images/data_spread_results.png}
\end{figure}

There is very little data represented by the high center and high right camera positions. This could lead to my initial observation that
higher positions are not as accurate being drowned out by the other data.
There is a surprising amount of hair\_in\_face data. This is because beards and fringes were both counted as hair\_in\_face, the main
dividing factor being the gender. Females with hair\_in\_face represents fringes or their hair down, while males with this attribute talk
about beards or the appearance of stubble.

The main observation is that overall accuracy is very good. With 58.7\% of the manually classified data being Alert and 63.9\% of the
models’ classifications also being Alert, this implies there is a strong consistency between the model and my own opinions. However,
you can also see from this that there is considerable discrepancy in my own classifications and that of the mode. But, despite this, the
model appears to make mistakes more in te middle of the range (when it is more unsure) suggesting that there are no extreme paths within
the model that have been aloud to spiral out of control (e.g. always classifying tired or always classifying alert). This is a good
sign as it suggests that the model is not overfitted to one classification, despite the slight bias in the training data.

\subsubsection{How \textbf{Darkness} affects the model’s Prediction:}
\begin{figure}[H]
    \centering
    % First Subfigure
    \begin{subfigure}{0.48\textwidth}
        \centering
        \includegraphics[width=\textwidth]{images/box_actual.png}
        \caption{Actual classification}
        \label{fig:box_actual}
    \end{subfigure}
    \hfill % This pushes the two images to the outer edges
    % Second Subfigure
    \begin{subfigure}{0.48\textwidth}
        \centering
        \includegraphics[width=\textwidth]{images/box_predicted.png}
        \caption{Predicted classification}
        \label{fig:box_predicted}
    \end{subfigure}

    \caption{Comparison of darkness variance between actual and predicted classifications.}
    \label{fig:darkness_comparison}
\end{figure}
With the position of the median shifting towards lower values, it could be safe to assume that the model appears to be become less
reliable in lower light conditions (especially when predicting tiredness). With the upper and lower quartiles also being smaller for
the predicted values, it suggests that the model is more likely to detect tiredness in lower light conditions (which is to be expected
as lower light implies earlier in the morning/later in the day – when someone is more likely to be tired).

\begin{figure}[H]
    \centering
    \includegraphics[width=0.8\textwidth]{images/correlations.png}
    \caption{Darkness effect}
\end{figure}

\begin{figure}[H]
    \centering
    \includegraphics[width=0.8\textwidth]{images/darknessxmodel_confidence.png}
    \caption{Darkness vs model confidence}
    \label{fig:darkness_confidence}
\end{figure}
With the R values (regression lines) in (figure \ref{fig:darkness_confidence}) being relitively low ($<0.3$) but
positive, suggests that there is a weak, but present, correlation between darkness and the model's confidence. This
shows that there is a subconcious tendancy within the model to predict a higher confidence for tiredness in lower
light (when the image is darker) conditions, despite the user presenting as \textbf{alert}.

The lower correlation ($R$) between darkness and confidence for tiredness suggests the model is more robust to
lighting changes. While \textbf{alert} predictions are more sensitive to light levels—potentially due to visibility
issues or over-exposure causing subjects to squint—the model remains more consistent in its confidence across
varying light conditions when identifying fatigue.


\paragraph{Observations on Data Density:} Within Figure \ref{fig:darkness_confidence}, certain vertical clusters of
data points are visible, particularly around a darkness level of $0.75$. Rather than representing a model bias, this
high density indicates:
    \begin{itemize}
    \item \textbf{Batch Consistency:} Large subsets of the testing data were likely captured under identical 
    environmental conditions, such as the same vehicle interior, specific time of day, or consistent camera sensor
    settings (or entire camera).
    \item \textbf{Discrete Metadata:} These "strips" represent the limitations of the darkness metadata , where
    multiple images are assigned the exact same darkness value despite slight visual variances.
\end{itemize}




\subsubsection{How well the mode performs over varying \textbf{genders}?} \label{sssec:gender section}
\begin{figure}[H]
    \centering
    \includegraphics[width=0.5\textwidth]{images/seive_gender.png}
    \caption{Sieve Diagram of Gender vs. Confidence}
    \label{fig:sieve_gender}
\end{figure}
\paragraph{Sieve Diagram:} This is used to show the relation between observed features and a 'null hypothesis' (how well the data matches
the expected number of occurences).  A statistically significant association is present ($\chi^2 = 261.66, p < 0.001$), indicating that
model confidence is not independent of gender:
\begin{itemize}
    \item Female Over-representation: A strong positive association (dark blue) exists for conficenve >=0.7.
    \item Male Under-representation: Conversely, the models is signifficantly less likely to assign a high confidence (far fewer times
    than statistically expected).
\end{itemize}

\begin{figure}[H]
    \centering
    % First Image
    \begin{subfigure}[b]{0.45\textwidth}
        \centering
        \includegraphics[width=\textwidth]{images/viola_gender(alert).png}
        \caption{Alert Confidence}
    \end{subfigure}
    \hfill % This adds the horizontal space to push them apart
    % Second Image
    \begin{subfigure}[b]{0.45\textwidth}
        \centering
        \includegraphics[width=\textwidth]{images/viola_gender(tired).png}
        \caption{Tired Confidence}
    \end{subfigure}
    \caption{Comparison of model confidence scores by gender.}
    \label{fig:gender_violins}
\end{figure}
\paragraph{Violin Plot Analysis:} While the sieve diagram identifies a frequency anomaly, the violin plots (Figure \ref{fig:gender_violins})
clarify the nature of this bias by visualizing the raw distribution of confidence scores
\begin{itemize}
    \item \textbf{Polarized Certainty:} \textbf{Famales} show a higher discriminatory power in terms of their classificaiton, with their
        \textbf{alert\_confidence} being skewed more towards the ceiling (1.0) and \textbf{tired\_confidence} being skewed more
        towards the floor (0.0). This suggests that the model is more certain of alertness for women.
    \item \textbf{Classification entropy:} For \textbf{males}, there is more of a centralised cluster (0.45--0.55). This indicates a higher
        level of uncertainty in the models predictions, resulting in more lower-conviction evaluations.
\end{itemize}
Overall, with women being only $12.99\%$ of the total tested dataset, they simply yield a higher level of predictive certainty. This is
likely due to higher clarity of training data which could have resulted in overfitting or lack of ambiguity in training data. This could
show that the model is more highly-optimised for female-centric characteristics.




\subsubsection{Effect of \textbf{accessories} on accuracy}
\begin{figure}[H]
    \centering
    % First Image
    \begin{subfigure}[b]{0.45\textwidth}
        \centering
        \includegraphics[width=\textwidth]{images/alert_vs_accessories.png}
        \caption{Alert Confidence}
    \end{subfigure}
    \hfill % This adds the horizontal space to push them apart
    % Second Image
    \begin{subfigure}[b]{0.45\textwidth}
        \centering
        \includegraphics[width=\textwidth]{images/tired_vs_accessories.png}
        \caption{Tired Confidence}
    \end{subfigure}
    \caption{Comparison of the effect Accessories have on confidence.}
    \label{fig:accessoried_box}
\end{figure}
In Figure \ref{fig:accessoried_box}, the presence of accessories (such as glasses, hats and masks) appeats to have a signigficant effect on the
models confidence. With both graphs showing a significant shift, you can see that the model tends more towards a \textbf{tired} 
classification when accessories are present.
\paragraph{Hypothesis Testing:}
With the sample size being $1767$ images (for the smaller classification -- with accessories), the data is sufficient to draw
statistically significant conclusions. With a mean difference of $0.142$ (14.2\%) in alert confidence between the two groups, and a
standard deviation of $0.25$, the calculated effect size (Cohen's d) is approximately $0.568$, which is considered a medium effect size.
This indicates that the presence of accessories has a moderate impact on the model's confidence in classifying alertness, with a
significant shift towards tired classifications when accessories are present. However, this does suggest that the model works flawlessly
with certain images where this is present, but also terribly on specific examples, which is not ideal for a model that is supposed to be 
versitile and work in a varitey of circumstances.

\paragraph{Further Considerations:} To further rienforce these findings, I have made use of the \textbf{violin} structure in Orange
\cite{orange} to illustrate this point. If this is true, the main body of the data should sit higher in accessory cases for tired
confidence, and lower for alert confidence. This is exactly what we see in the figure below:
\begin{figure}[H]
    \begin{subfigure}[b]{0.45\textwidth}
        \centering
        \includegraphics[width=\textwidth]{images/alert_vs_accessories(violin).png}
        \caption{Alert Confidence}
    \end{subfigure}
    \hfill
    \begin{subfigure}[b]{0.45\textwidth}
        \centering
        \includegraphics[width=\textwidth]{images/tired_vs_accessories(violin).png}
        \caption{Tired Confidence}
    \end{subfigure}
\end{figure}




\subsubsection{How does a user having their \textbf{hair in their face} affect the model?}
\begin{figure}[H]
    \centering
    \includegraphics[width=0.4\textwidth]{images/hair_in_face (fairness).png}
    \caption{Dataset Bias Indication}
\end{figure}
\begin{figure}[H]
    \centering
    \includegraphics[width=\textwidth]{images/fairness_nodes.png}
    \caption{Node structure for fairness testing}
    \vspace{2pt}
    \small{\textbf{Note:} The warning sign in the Dataset Bias widget refers to missing values within specific feature columns.
    These issues are omitted from the Disparate Impact calculation to maintain the validity.}
\end{figure}
With a \textbf{Disparate Impact} value of $0.999$, there is very little evidence to suggest that the model is unfairly biased to either
classificaiton when the subject has their \textbf{hair\_in\_face}. The \textbf{Parity Difference} represents the difference in how the
model treats the different groups. In this case, there is $-0.0$ difference, which implies there is a tiny proportion (uncountably small)
tat points to one classificaiton over another. However, this isso small that it could simply represent a single instance which was
misslabeled in the entire dataset. Thus implying that the model is 'blind' to the effects of this classificaiton.
\begin{figure}[H]
    \centering
    \includegraphics[width=\textwidth]{images/hair_in_face (alert-scatter).png}
\end{figure}
This is re-enforced with the scatter plot where you can see there is no correlation between this and the model's confidence. To account
for the linear and unclear nature of duel classifications, the 'Jitter' of the data has been scaled to the maximum to form a maximum
disparity between data points to allow for better visualiisation.

\paragraph{Hair positioning:}
Within this classification of \textbf{hair\_in\_face} there are two different representations:
\begin{itemize}
    \item \textbf{Beards} \- when on a male user
    \item \textbf{hair down} \- when matched with a female participant
\end{itemize}
While you can have cross-overs in this representation (e.g. men with long hair that can over their face), throught all of the datasets,
there was no-one who countered this assumption, so it was easier to classify them as the same (with the distinguishing factor being
gender).
Since gender bias was already determined in \ref{sssec:gender section}, that women yield a greater predictive certainty,
it would be wise to compair the combined effect of these characteristict. Leading to the question:
\\\texttt{Is the greater certainty of a classification a result
of a greater presence of hair?}
\begin{figure}[H]
     \centering
     \begin{subfigure}[b]{0.3\textwidth}
         \centering
         \includegraphics[width=\textwidth]{images/gender_hair_box.png}
         \caption{Model confidence across subgroups}
         \label{fig:gender_hair_box}
     \end{subfigure}
     \hfill
     \begin{subfigure}[b]{0.6\textwidth}
         \centering
         \includegraphics[width=\textwidth]{images/gender_hair_correct.png}
         \caption{Classification accuracy distribution}
         \label{fig:gender_hair_correct}
     \end{subfigure}
     \caption{Comparison of model confidence and accuracy for Gender and Hair combinations}
     \label{fig:combined_gender_hair_analysis}
\end{figure}
While the overall data suggests that there is no effect in overall accuracy as a result of a user haivng their \textbf{hair\_in\_face}, figure
\ref{fig:combined_gender_hair_analysis} suggests this is not the case regarding its placement. For male users (where it primarilly represents beards), the presence
of hair correlates to a decrease in models accuracy, suggesting facial hear acts as \textbf{noise} for feature extraction. Conversely, for female participants
(representing hair down), the model shows an increase in accuracy and certainty. This suggests that while beards obscure crutial features, long hair inadvertantly
helps to act as a secondary framing feature, helping the model anchor the face.




\subsubsection{Camera position effect on accuracy}
\begin{figure}[H]
    \centering
    \includegraphics[width=0.8\textwidth]{images/tirednesxcamera_pos.png}
    \caption{Camera Position vs Model Confidence of Classifications}
\end{figure}
\noindent
\begin{minipage}{\linewidth}
Which represents the positions as follows (as you look towards the driver):\\[4pt]
\begin{center}
    \begin{tabular}{|c|c|c|}
        \hline
        0 & 1 & 2 \\ \hline
        3 & 4 & 5 \\ \hline
        6 & 7 & 8 \\ \hline
    \end{tabular}
\end{center}
\end{minipage}

As you can see, the position of the camera has significant effect on accuracy. With higher positions being more
likely to be classified as tired (potentially as a result of eyelids being considered more closed with respect to
the rest of the eye), while lower positions are more likely to be classified as alert. This is a significant issue
as it suggests that the model is not as versatile as I would have hoped, and that the point I initially made about
hardware and vehicle versitility is not met.
\begin{figure}[H]
    \centering
    \includegraphics[width=0.8\textwidth]{images/camera_pos_node.png}
    \caption{Camera Position separation}
\end{figure}

\begin{figure}[H]
     \centering
     \begin{subfigure}[b]{0.48\textwidth}
         \centering
         \includegraphics[width=\textwidth]{images/true_tired_cam.png}
         \caption{Alert prediction in Tired cases}
     \end{subfigure}
     \hfill
     \begin{subfigure}[b]{0.48\textwidth}
         \centering
         \includegraphics[width=\textwidth]{images/true_alert_cam.png}
         \caption{Alert prediction in Alert cases}
     \end{subfigure}
     \caption{Positive predictions across the camera locations}
     \label{fig:camera_distributions}
\end{figure}
In the Figure above, a distinct correlation between camera position and the model's confidence on True values is evident. The results
demonstrate that the model achieves significantly higher accuracy when the camera is in the Middle position. This aligns with
expectations, as the training dataset consists primarily of centered, front-facing imagery.
\paragraph{Hypothesis Testing:}
Even in the smallest sample group (Top: $N=587$), the sample size exceeds the minimum requirement of 384 units for statistical
significance. We established that a discrepancy greater than 5\% between groups would constitute substantial evidence of a "Position
Effect."The observed delta between the Middle position ($0.46$) and the average of the peripheral positions ($\approx 0.37$) is $0.09$
(a 9\% difference). Because this 9\% gap exceeds our 5\% threshold—and given that the 95\% confidence intervals for the Middle group do
not overlap with the others—there is significant evidence to conclude that camera position has a substantial effect on model performance,
with the central position yielding the highest accuracy.
High confidence however is only valuable if is correlates to correct classificaitons.
\begin{figure}[H]
    \centering
    \includegraphics[width=0.7\textwidth]{images/correct_cam.png}
    \caption{proportion bars for correctness over camera positions}
    \label{fig:correct_cam_accuracy}
\end{figure}

As illustrated in Figure \ref{fig:correct_cam_accuracy}, the model is undoubtadly more accurate (and as shown previously) more certain of
a specific classification when the camera is centre mounted, compared to other positions.



\subsection{Overall Accuracy} 
\begin{figure}[H]
    \centering
    \includegraphics[width=0.5\textwidth]{images/overall_accuracy.png}
\end{figure}
Using the Python node, I counted through every line and counted the differences between the model’s prediction and the actual
classification. With $50/50$ being the average split for simple guesses, independent of any image, a $77\%$ classification is great for a
variety of circumstances. Considering that this was all trained off of one dataset, which consisted of a fixed camera angle, same vehicle,
same lighting conditions and a limited number of participants with little variation in terms of accessory features, I am more than happy
with this final result.
\subsubsection{Distribution of Results:}
\begin{figure}[H]
    \centering
    \includegraphics[width=0.5\textwidth]{images/true-pos-true-neg.png}
\end{figure}
The distribution above illustrates the model's predictive performance across the entire dataset, specifically focusing on the 5,675 complete
rows of data. Analysis of the classification errors indicates that the model is more inclined toward False Negative predictions regarding
alert states (mistakenly identifying tiredness). From a safety-critical perspective, this bias is preferable; it is more advantageous for the
system to err on the side of caution by providing a 'false alarm' than to fail to detect a genuine instance of driver fatigue, which would pose
a significantly higher risk.

\subsubsection{Confidence in False Cases}
\begin{figure}[H]
    \centering
    \includegraphics[width=0.4\textwidth]{images/conf_diff_wrong.png}
    \caption{Distribution of Peak Classification Confidence vs Accuracy}
    \label{fig:confidence vs correction}
\end{figure}
The model demonstrates a desirabel \textbf{"self-aware"} error profile, where the predictive certainty is reduced in cases where the model is wrong. With
Figure \ref{fig:confidence vs correction} showing that incorrect classificaitons exhibit a significantly greater density cluster near the 0.5 mark (border between 2
calsses), with correct classifications haviung a \textbf{higher median confidence}. This is statistically supported by the moderate correlation between
confidence levels and classificairtion outcomes:
\begin{figure}[H]
    \centering
    \includegraphics[width=0.5\textwidth]{images/error_rate_correlation.png}
\end{figure}






\subsubsection{Summary of Attribute Effects and Ideal Environment}
To quantify the impact of each variable, the model's performance was measured against the overall baseline accuracy of 77\%. The "Delta" represents the percentage point
shift in accuracy for this atrribute.

\begin{landscape}
%\begin{center}
\renewcommand{\arraystretch}{1.3} % Adds vertical padding for better readability
\begin{xltabular}{\linewidth}{@{} l c c X @{}}
\caption{Summary of Attribute Effects on Model Accuracy} \label{tab:attr_summary} \\
\toprule
    \textbf{Attribute} & \textbf{Acc. Overall} & \textbf{Delta} & \textbf{Key Insight \& Ideal Environment} \\
\midrule
\endfirsthead

\toprule
    \textbf{Attribute} & \textbf{Acc. Overall} & \textbf{Delta} & \textbf{Key Insight \& Ideal Environment} \\
\midrule
\endhead

Accessories (None)    & $83.9\%$  & $+6.6\%$  & Presence of eyewear/masks shifts bias toward \textbf{Alert} ($91.16\%$). Sub-optimal for safety-critical environments where false negatives are risky. \\ \addlinespace
Accessories (Present) & $62.76\%$ & $-14.5\%$ & Superior at predicting \textbf{Tiredness} ($83.54\%$). From a safety standpoint, this "erring on the side of caution" is preferable for fatigue detection. \\ \midrule

Gender (Female)       & $82.2\%*$ & $+4.9\%$  & High alert accuracy; however, "Tired" cases show significantly lower parity ($44.7\%$), suggesting potential training set imbalance. \\ \addlinespace
Gender (Male)         & $76.6\%$  & $-0.7\%$  & Highly balanced across classes ($\pm 4.4\%$); significantly less prone to extreme classification bias than female subjects. \\ \midrule

Camera (Top)          & $59.6\%*$ & $-17.7\%$ & Over-weighted toward \textbf{Tired} due to downward eyelid perspective; useful specifically for prioritizing safety-critical "false alarms." \\ \addlinespace
Camera (Middle)       & $88.3\%$  & $+11.0\%$ & Optimal for predictive confidence; aligns with the primary training distribution (centered, front-facing imagery). \\ \addlinespace
Camera (Bottom)       & $75.7\%$  & $-1.6\%$  & Most stable and versatile across both classifications with a minimal recorded bias of only $0.16\%$. \\ \midrule

Hair in Face (None)   & $86.1\%$  & $+8.8\%$  & High predictive power but exhibits greater disparity between classes, with a $94.8\%$ bias toward \textbf{Alert} classifications. \\ \addlinespace
Hair in Face (Present)& $73.1\%$  & $-4.2\%$  & Context-dependent: Facial hair acts as noise for male subjects, while long hair assists in framing for female participants. Minimal accuracy delta of $3.2\%$. \\

\bottomrule
\multicolumn{4}{l}{\small{*Note: High accuracy in these specific subsets is often a result of class-weighted bias rather than general model power.}}
\end{xltabular}%
%\end{center}%
\end{landscape}
\nopagebreak

\paragraph{Darkness Distribution}
\begin{figure}[H]
    \centering
    \includegraphics[width=0.7\textwidth]{images/brightnerss_correction_relation.png}
\end{figure}

There is little correlation between the brightness and the corrrectness of the model. While it appears that in birgh conditions the model works best, but this is not the case as the $\leq0.9$
only represents $0.04\%$ of the data and so is not representative of the dataset as a whole. The same can be said for $<0.4$ where it represents only $\approx2\%$. As a result, the ideal range is $0.5\to0.8$ as this has a high succeess rate and has enough data to
support this conclusion, as each subsequent section represents around $25\%$ of the data each.

\subsection{Ideal Use Case / Environment:}
Based on the aggregated data, the "Perfect Environment" for this model is a user with long hair, no facial accessories, situated in a vehicle with a center-bottom
mounted camera, and operating in a darkness index between $0.5$ and $0.8$. Deviations from these parameters—specifically the introduction of eyewear or side-mounted cameras—require
a lower confidence threshold for alerts to compensate for the model's inherent biases. Leading to an increased chance of misclassificaiotn.



\subsection{User Feedback:}
    \TODO{get some people to use it, take some pictures of placements and other things - talk about what they liked/didn't like}

\subsubsection{Areas of Success}
    \TODO{areas of success}
\subsubsection{Areas of Failure}
    \TODO{what could I have done better if done again, or areas that were difficult that really stand out}

\subsection{Additional Remarks and Development ideas}
This section will outline any ideas or areas that I beleive are worth commenting on; with many of these being future considerations given
more time.